\section{Uniform spectral and entropy estimates}
\label{rcut:sec:entropy}

The entropy estimate needed for the independent cutoff argument rests on a
spectral gap that survives both changes of external field and reductions of
the interaction strength.  We prove this gap by summing a signed two-spin
inequality, then pass to entropy through Gaussian observations.  The observation
posterior changes both the field and the interaction, which is why both
uniformities are needed.  Throughout this section, the inverse temperature
$\beta<1/2$ is fixed.

\subsection{Gaussian estimates and a good disorder event}

We first collect the Gaussian tools used to control the operator norm, the
row-square sums, and the largest entry of the interaction matrix.  These three
controls will make the two-spin error summable on a single disorder event.

\begin{lemma}[Gaussian concentration and comparison]
\label{rcut:lem:gaussian-facts}
Let $Z$ be a vector of independent standard real Gaussian variables.
If $F$ is $D$-Lipschitz in Euclidean distance, then
\begin{equation}
 \E\exp\bigl(s(F(Z)-\E F(Z))\bigr)
 \le \exp\left(\frac{\pi^2D^2s^2}{8}\right),
 \qquad s\in\R.
 \label{rcut:eq:gaussian-mgf}
\end{equation}
In particular, for $D>0$ and $r>0$,
\[
 \P\bigl(|F(Z)-\E F(Z)|\ge r\bigr)
 \le 2\exp\left(-\frac{2r^2}{\pi^2D^2}\right).
\]
If $(X_a)_{a\in I}$ and $(Y_a)_{a\in I}$ are centered Gaussian processes on
the same finite index set and
\[
 \E(X_a-X_b)^2\le \E(Y_a-Y_b)^2
 \quad\text{for all }a,b\in I,
\]
then $\E\max_a X_a\le\E\max_a Y_a$.
The comparison also holds for continuous processes on a compact metric index
set when their suprema are integrable.
\end{lemma}

\begin{proof}
For the concentration estimate, we use the Maurey--Pisier Gaussian
rotation and double-Jensen argument; see Giraud~\cite[Exercise~1.6.7,
equation~(1.10), parts~A--B]{Giraud2021}. Rotation between two independent
Gaussian vectors expresses the increment as an integral of directional
derivatives.
First take $F$ smooth with $\|\nabla F\|\le D$.
Let $Z'$ be an independent copy of $Z$ and define the rotation by
\[
 Z_\theta=Z\cos\theta+Z'\sin\theta,
 \qquad V_\theta=-Z\sin\theta+Z'\cos\theta,
 \qquad 0\le\theta\le\pi/2.
\]
For a fixed angle, $Z_\theta$ and $V_\theta$ are independent standard
Gaussian vectors.  This independence lets us evaluate the exponential
moment after applying Jensen's inequality, first in $Z'$ and then with
respect to the uniform measure on $[0,\pi/2]$:
\begin{align*}
 \E e^{s(F(Z)-\E F(Z))}
 &\le \E e^{s(F(Z)-F(Z'))}\\
 &\le \frac2\pi\int_0^{\pi/2}
     \E\exp\left(-\frac{\pi s}{2}
             \nabla F(Z_\theta)\cdot V_\theta\right)\,d\theta
 \le e^{\pi^2D^2s^2/8}.
\end{align*}
To pass to a general Lipschitz $F$, smooth by convolution and use dominated
convergence.  Linear growth gives the Gaussian exponential integrability
required for this passage.  Optimization of the exponential Markov bound
then gives the tail estimate.  The case $D=0$ follows directly because
the function is constant.

For the comparison statement, we interpolate independent realizations of
the two processes and show that the expected maximum increases along the
interpolation.  For $\lambda>0$, replace the maximum by the smooth function
\[
 \Phi_\lambda(z)=\lambda^{-1}\log\sum_{a\in I}e^{\lambda z_a}.
\]
Gaussian integration by parts gives
\[
 \frac{d}{dt}\E\Phi_\lambda
       \bigl(\sqrt{1-t}\,X+\sqrt t\,Y\bigr)
 =\frac12\sum_{a,b}B_{ab}\,\E(\partial_{ab}\Phi_\lambda),
 \qquad B=\Cov(Y)-\Cov(X).
\]
The Hessian has zero row sums and nonpositive off-diagonal entries.
Consequently the derivative can be expressed using only increment
variances.  With $\Delta_{ab}=B_{aa}+B_{bb}-2B_{ab}$ denoting their
difference, the right side becomes
\[
 -\frac14\sum_{a,b}\Delta_{ab}\,\E(\partial_{ab}\Phi_\lambda)\ge0.
\]
The limit $\lambda\to\infty$ gives the finite-set comparison.  Exhausting
a dense set by finite subsets and using continuity and dominated convergence
gives the compact-set version.  In the applications below, every such
process is a linear function of a finite Gaussian vector, which ensures
integrability of its supremum.
\end{proof}

\begin{proposition}[Good disorder]
\label{rcut:prop:good-disorder}
There are deterministic numbers $\delta_n\downarrow0$ and events
$\mathcal E_n=\mathcal E_n(\beta)$ with $\P(\mathcal E_n)\to1$ such that
on $\mathcal E_n$,
\begin{equation}
 \|J\|\le2\beta+\delta_n,
 \qquad \max_i\sum_jJ_{ij}^2\le\beta^2+\delta_n,
 \qquad \max_{i,j}|J_{ij}|\le\delta_n.
 \label{rcut:eq:good-disorder}
\end{equation}
Consequently, on the same events,
\[
 \max_i\sum_j|J_{ij}|^3\longrightarrow0.
\]
One may take $\delta_n=n^{-1/4}$.
\end{proposition}

\begin{proof}
We prove the operator-norm estimate first and then control the row sums
and individual entries.  The case $\beta=0$ is deterministic.  For
$\beta>0$, add independent diagonal entries of variance $2\beta^2/n$ to
$J$.  The resulting centered Gaussian orthogonal ensemble (GOE) matrix
$W$ has a covariance that permits a direct Gaussian comparison: for unit
vectors $p,q$,
\[
 \E[(p^{\mathsf T}Wp)(q^{\mathsf T}Wq)]
 =\frac{2\beta^2}{n}(p^{\mathsf T}q)^2.
\]
With $\rho=p^{\mathsf T}q$, the increment variance of $p^{\mathsf T}Wp$ is
$4\beta^2(1-\rho^2)/n$.  Bound it by $8\beta^2(1-\rho)/n$ and compare
with the linear process $(2\beta/\sqrt n)p^{\mathsf T}z$, where $z$ is
a standard Gaussian vector.  Lemma~\ref{rcut:lem:gaussian-facts} yields
\[
 \E\lambda_{\max}(W)
 \le\frac{2\beta}{\sqrt n}\E\|z\|\le2\beta.
\]
In the independent standard Gaussian coordinates defining $W$, the
quadratic form $p^{\mathsf T}Wp$ is Lipschitz with constant
$\sqrt2\beta/\sqrt n$, uniformly over unit vectors.  The supremum inherits
this constant.  Concentration for $\lambda_{\max}(W)$ and
$\lambda_{\max}(-W)$ therefore bounds $\|W\|$ by $2\beta+\delta_n/2$
with probability tending to one.  The largest absolute diagonal entry is
at most $\delta_n/2$ with probability tending to one as well.  Removing
that diagonal proves the first estimate in
Equation~\eqref{rcut:eq:good-disorder}.

For a standard Gaussian $g$,
\[
 \E e^{\theta(g^2-1)}=e^{-\theta}(1-2\theta)^{-1/2}
 \le e^{C\theta^2},\qquad |\theta|\le1/4.
\]
Apply the exponential Markov bound to the $n-1$ independent squares in
one row.  The probability that their weighted sum exceeds
$\beta^2+\delta_n$ is at most $\exp(-c_\beta n\delta_n^2)$ for large $n$.
A union bound over the rows gives the second estimate; Gaussian tails
and a union bound over entries give the third.  The choice
$n\delta_n^2=\sqrt n$ makes all the failure probabilities vanish.
The cubic row sums needed for the two-spin remainder are now controlled
by the last two estimates together:
\[
 \max_i\sum_j|J_{ij}|^3
 \le \left(\max_{i,j}|J_{ij}|\right)
       \left(\max_i\sum_jJ_{ij}^2\right)
 \le\delta_n(\beta^2+\delta_n).
\]
\end{proof}

\subsection{A two-spin estimate and a uniform gap}

We next prove an inequality for a pair of sites and sum it over the
set of site pairs to obtain the gap.  The pair estimate must remain
uniform when a one-site conditional variance is small.  For an Ising
law $\nu$ with symmetric, zero-diagonal interaction matrix $A$ and
external field $\theta\in\R^n$, write
\[
 \nu(x)\propto\exp\left(\frac12x^{\mathsf T}Ax+\theta\cdot x\right),
 \qquad m_i=\E_\nu[X_i\mid X_{-i}],\qquad v_i=1-m_i^2.
\]
All notation in this subsection refers to this law.  To express the
one-site energies as squared projection norms, let $E_i$ be conditional
expectation given the spins outside $i$, and put
\[
 \mathsf D_i=I-E_i,\qquad \mathsf D_i f=(x_i-m_i)d_i f.
\]
These are orthogonal projections in $L^2(\nu)$, and
\[
 \nu[(\mathsf D_i f)^2]=\nu[v_i(d_i f)^2],
 \qquad -L_\nu=\sum_i\mathsf D_i,
\]
where $L_\nu$ is the heat-bath generator with rate one at every site.

The signed two-spin estimate below follows the gap method of
Wang~\cite[Theorem~1.2 and Lemma~3.4]{Wang2026}.  Keeping the sign of the
leading interaction term allows the later summation to use an operator
norm.  We give the algebra and the remainder estimate that make this
summation uniform in the external fields and the interaction scale.

\begin{lemma}[Uniform two-spin inequality]
\label{rcut:lem:two-spin}
There are absolute constants $t_0>0$ and $C_0<\infty$ with the following
property.  For any finite-field Ising law $\nu$, any distinct $i,j$ with
$t=\tanh A_{ij}$ satisfying $|t|\le t_0$, and any real function $f$,
\begin{align}
 \nu[\mathsf D_i f\mathsf D_j f]
 &\ge-\nu\bigl[(t+2t^2m_i m_j)v_i v_jd_i f d_j f\bigr]\notag\\
 &\hspace{1cm}-C_0|t|^3\nu\bigl[(\mathsf D_i f)^2+(\mathsf D_j f)^2\bigr].
 \label{rcut:eq:two-spin}
\end{align}
The constants are independent of all external fields and of the dimension.
For example, $t_0=1/100$ and $C_0=22000$ are valid.
\end{lemma}

\begin{proof}
The proof has three steps: factor out the two conditional variances,
identify the nonnegative part of the expansion, and absorb the cubic
remainder into the one-site energies.  Condition on all spins outside
$i,j$, and denote the remaining spins by $x,y$.  Their probabilities are
\[
 \frac{(1+ax)(1+by)(1+txy)}{4(1+abt)},
 \qquad |a|,|b|<1,
\]
where $a,b$ are the hyperbolic tangents of the external fields left after
conditioning.  Set $Z_0=1+abt$.  The corresponding conditional means
and variances take the form
\[
 m_i(y)=\frac{a+ty}{1+aty},\qquad
 v_i(y)=\frac{(1-a^2)(1-t^2)}{(1+aty)^2},
\]
and the analogous expressions with $a,x$ and $b,y$ interchanged.
There are real numbers $p,q,r$ such that
\[
 d_i f=p+r(y-b),\qquad d_j f=q+r(x-a).
\]
Indeed, expand the restriction of $f$ in the basis $1,x,y,xy$ and
denote the coefficient of $xy$ by $r$.  This isolates the common
second-difference coefficient from the two first differences.

Let $\E_0$ denote uniform expectation over the four pairs $(x,y)$.
We now combine the mixed projection energy with the correction term in
the claimed inequality.  The identities
\[
 x-m_i(y)=\frac{(x-a)(1-txy)}{1+aty},
 \qquad (1+ax)(x-a)=x(1-a^2)
\]
and their counterparts for $y$ give the exact factorization
\begin{align*}
 &\nu[\mathsf D_i f\mathsf D_j f\mid X_{-\{i,j\}}]
 +\nu[(t+2t^2m_i m_j)v_i v_jd_i f d_j f\mid X_{-\{i,j\}}]
   =P_0Q,\\
 &P_0=\frac{(1-a^2)(1-b^2)(1-t^2)}{Z_0},\qquad
 Q=\E_0\bigl([p+r(y-b)][q+r(x-a)]C(t;x,y)\bigr),
\end{align*}
where
\begin{align*}
 C(t;x,y)
 ={}&\frac{xy-t}{(1+aty)(1+btx)}\\
 &+(t+2t^2m_i(y)m_j(x))
  \frac{(1-t^2)(1+ax)(1+by)(1+txy)}
       {(1+aty)^2(1+btx)^2}.
\end{align*}
Only the expansion through second order will contribute to the leading
term.  Expanding at $t=0$ and reducing by $x^2=y^2=1$ gives
\begin{equation}
 C(t;x,y)
 =xy\left[1+ab t+(1-a^2-b^2+2a^2b^2)t^2\right]
   +R(t;x,y).
 \label{rcut:eq:two-spin-expansion}
\end{equation}
The cancellation can be checked directly: before collecting terms, the
coefficient of $t^2$ is
\[
 d+xy(d^2-abxy)+(1+ax)(1+by)(xy-2d+2ab),
 \qquad d=ay+bx.
\]
Collecting terms gives the coefficient in
Equation~\eqref{rcut:eq:two-spin-expansion}, with the coefficients of
$1,x,y$ cancelling.  Thus the displayed expansion contributes through
the following single identity:
\[
 \E_0\bigl(xy[p+r(y-b)][q+r(x-a)]\bigr)=r^2.
\]

It remains to control the remainder without excluding nearly
deterministic spins; we therefore allow the limiting values $|a|=|b|=1$
in this estimate.  On the complex disk $|t|\le1/2$, the denominators
$1+aty$ and $1+btx$ have modulus at least $1/2$, while
$|m_i(y)|,|m_j(x)|\le3$.  The first term of $C$ is bounded in modulus
by $6$.  In the second term, $t+2t^2m_i m_j$ is bounded by $5$;
including its remaining factors gives a bound of $600$.
Hence $|C(t;x,y)|\le606$ throughout the disk, and Cauchy's coefficient
estimate yields
\[
 |R(t;x,y)|\le606\sum_{k\ge3}(2|t|)^k
 \le9696|t|^3,\qquad |t|\le1/4.
\]
Since
\[
 |[p+r(y-b)][q+r(x-a)]|\le p^2+q^2+8r^2
\]
and $1-a^2-b^2+2a^2b^2=a^2b^2+(1-a^2)(1-b^2)\ge0$, we obtain
\[
 Q\ge\bigl(1-|t|-8\cdot9696|t|^3\bigr)r^2
          -9696|t|^3(p^2+q^2).
\]
For $|t|\le1/100$, the coefficient of $r^2$ exceeds $1/2$.

The leading term is now nonnegative.  To finish, the negative cubic term
must be charged to the conditional energies while retaining the variance
factors.  Write $A_0=1-a^2$ and $B_0=1-b^2$.  Summing over $x$ first
gives the exact energy formula
\[
 \nu[v_i(d_i f)^2\mid X_{-\{i,j\}}]
 =\frac{A_0(1-t^2)}{2Z_0}
   \sum_{y=\pm1}\frac{1+by}{1+aty}[p+r(y-b)]^2.
\]
For $|t|\le1/4$, the denominators lie between $3/4$ and $5/4$, so
\begin{align*}
 \nu[v_i(d_i f)^2\mid X_{-\{i,j\}}]
 &\ge\frac35 A_0(p^2+B_0r^2),\\
 \nu[v_j(d_j f)^2\mid X_{-\{i,j\}}]
 &\ge\frac35 B_0(q^2+A_0r^2),\qquad
 P_0\le\frac43A_0B_0.
\end{align*}
These bounds make $P_0p^2$ and $P_0q^2$ at most $20/9$ times their
respective conditional energies.  Multiply the estimate for $Q$ by
$P_0$ and use $9696\cdot20/9<22000$ to obtain the conditional
inequality.  Averaging over the spins on which we conditioned gives
Equation~\eqref{rcut:eq:two-spin}.  The estimates are homogeneous in
$p,q,r$, and throughout the argument neither these coefficients nor
the external fields were bounded.
\end{proof}

\begin{proposition}[Gap and covariance under scaled interactions and fields]
\label{rcut:prop:uniform-gap}
For all sufficiently large $n$, uniformly over $J\in\mathcal E_n$,
$s\in[0,1]$, and $\theta\in\R^n$, let
\[
 \nu_{s,\theta}(x)\propto
 \exp\left(\frac s2x^{\mathsf T}Jx+\theta\cdot x\right).
\]
Its rate-one-per-site heat-bath generator has spectral gap at least
\[
 \kappa_\beta=\frac{1-2\beta}{2}>0.
\]
Moreover,
\begin{equation}
 \Var_{\nu_{s,\theta}}(f)
 \le\kappa_\beta^{-1}\sum_i
       \nu_{s,\theta}[v_i(d_i f)^2],
 \qquad
 \|\Cov_{\nu_{s,\theta}}(X)\|\le\kappa_\beta^{-1}.
 \label{rcut:eq:uniform-gap-covariance}
\end{equation}
\end{proposition}

\begin{proof}
We use the pair estimate to compare the squared generator with its
Dirichlet form.  Fix $\nu=\nu_{s,\theta}$ and define the zero-diagonal
matrix by $T_{ij}=\tanh(sJ_{ij})$.  On $\mathcal E_n$, all entries are
small enough for Lemma~\ref{rcut:lem:two-spin} to apply simultaneously.
At each configuration put $z_i=v_i d_i f$ and $z'_i=m_i z_i$.
Summation over ordered pairs gives
\begin{align*}
 \left\|\sum_i\mathsf D_i f\right\|_{L^2(\nu)}^2
 \ge{}&\sum_i\|\mathsf D_i f\|_{L^2(\nu)}^2
       -\nu[z^{\mathsf T}Tz+2(z')^{\mathsf T}(T\circ T)z']\\
 &-2C_0\left(\max_i\sum_j|T_{ij}|^3\right)
               \sum_i\|\mathsf D_i f\|_{L^2(\nu)}^2,
\end{align*}
where $\circ$ denotes entrywise multiplication.  The field-dependent
factors remain inside the vectors, so the interaction matrices can be
bounded by their operator norms.  With
$M=\max(\|T\|,2\|T\circ T\|)$, the sum of the two quadratic forms is
at most
\[
 M\sum_i v_i^2(1+m_i^2)(d_i f)^2
 \le M\sum_i v_i(d_i f)^2,
\]
because $v_i^2(1+m_i^2)=v_i(1-m_i^4)\le v_i$.

We next control these operator norms using only the disorder event.
The Taylor estimate $|\tanh u-u|\le C|u|^3$, together with the
maximum absolute row-sum bound for symmetric matrices, gives
\[
 \|T-sJ\|\le C\max_i\sum_j|J_{ij}|^3=o(1).
\]
Also $\|T\circ T\|\le\max_i\sum_jT_{ij}^2\le\beta^2+o(1)$.
The error bounds depend only on the deterministic tolerances defining
the event; they are uniform over $s,\theta$ and $J\in\mathcal E_n$.
The fixed assumption $\beta<1/2$ therefore gives
\[
 M\le\max(2\beta,2\beta^2)+o(1)=2\beta+o(1).
\]
Thus, for $A=-L_\nu=\sum_i\mathsf D_i$ and all sufficiently large $n$,
\[
 \|Af\|_2^2\ge\kappa_\beta\langle f,Af\rangle_\nu.
\]
To read this as a spectral gap, use that $A$ is self-adjoint and
nonnegative.  Every one-site conditional law has full support, so
$\mathsf D_i f=0$ for all $i$ forces $f$ to be constant along every
edge of the cube.  Thus the kernel consists exactly of constants.
Applied to a nonconstant eigenfunction, the last inequality bounds its
eigenvalue below by $\kappa_\beta$.  This proves the gap and variance
assertions.  The covariance conclusion is the linear-function case
$f(x)=c^{\mathsf T}x$:
\[
 c^{\mathsf T}\Cov_\nu(X)c
 \le\kappa_\beta^{-1}\sum_i\nu[v_i]c_i^2
 \le\kappa_\beta^{-1}\|c\|^2.
\]
\end{proof}

\subsection{Entropy through Gaussian observations}

We now turn the uniform covariance estimate into entropy dissipation.
For a probability measure $\nu$ and a positive function $f$, the
homogeneous entropy functional is
\[
 \Ent_\nu(f)=\nu[f\log f]-\nu[f]\log\nu[f].
\]
The convention $0\log0=0$ extends it to nonnegative functions.
Proposition~\ref{rcut:prop:entropy} imposes no normalization on $f$.
The proof uses the Gaussian observation entropy factorization principle of
Chen and Eldan~\cite[Theorem~49 in arXiv:2203.04163v2]{CE2025}.  We include the argument to keep
the interaction scaling and the rate-one-per-site clock explicit.

\begin{proposition}[Entropy inequality and reduction from point starts]
\label{rcut:prop:entropy}
There are constants $C_1,C_2<\infty$, depending only on the fixed
$\beta<1/2$, such that for all sufficiently large $n$ and
$J\in\mathcal E_n$, the zero-field Gibbs measure satisfies
\begin{equation}
 \Ent_\mu(f)\le C_1\mu\bigl[f(-L)\log f\bigr],
 \qquad f>0.
 \label{rcut:eq:entropy-inequality}
\end{equation}
For every initial point $x$ and every $t\ge0$, if
$\rho_{t,x}=S_t(x,\cdot)$, then
\begin{equation}
 \Ent_\mu\left(\frac{d\rho_{t,x}}{d\mu}\right)
 \le C_2n e^{-t/C_1}.
 \label{rcut:eq:entropy-decay}
\end{equation}
In particular, the continuous-time worst-start mixing time at every fixed
positive accuracy is $O_\beta(\log n)$, uniformly on these events.
\end{proposition}

\begin{proof}
We interpolate from the Gibbs measure to product posteriors, compare
their expected entropy with the initial entropy, and then remove the
observation from the one-site terms.  For large $n$,
Equation~\eqref{rcut:eq:good-disorder} ensures $\|J\|<1$.
Choose $K=J+2I$; then $K$ is positive definite and satisfies $\|K\|<3$.
Keep it fixed throughout the argument for the realized disorder.
Given $X\sim\mu$ and an independent standard Gaussian vector $z$, use
the observation, for $0<u\le1$,
\[
 Y_u=uK^{1/2}X+\sqrt u\,z.
\]
The posterior law $\nu_{u,y}=\operatorname{Law}(X\mid Y_u=y)$ is
\begin{equation}
 \nu_{u,y}(x)\propto\mu(x)
    \exp\left(y^{\mathsf T}K^{1/2}x-\frac u2x^{\mathsf T}Kx\right).
 \label{rcut:eq:gaussian-posterior}
\end{equation}
Discarding the diagonal term, which is constant on the cube, leaves
interaction matrix $(1-u)J$.  The observation creates an external field,
and a further linear tilt changes only that field.  Thus both uniformities
in Proposition~\ref{rcut:prop:uniform-gap} apply: every posterior and
every such tilt has covariance norm at most $C=\kappa_\beta^{-1}$.
At $u=1$, the interaction vanishes and the posterior is a product measure.

\medskip
\noindent\emph{Derivative of the expected posterior entropy.}
To quantify the entropy lost to the observation, set
$H(u)=\E[\Ent_{\nu_{u,Y_u}}(f)]$.  The derivative identity we need is
\begin{equation}
 H'(u)=-\frac12\E\left[
       \frac{\|K^{1/2}\Cov_{\nu_{u,Y_u}}(X,f)\|^2}
            {\nu_{u,Y_u}[f]}\right].
 \label{rcut:eq:posterior-entropy-derivative}
\end{equation}
The covariance with $f$ here is a vector.  We verify the identity using
the finite Gaussian mixture defining the observation.  Let $\phi_u$ be
the centered Gaussian density with covariance $uI$, and define
\begin{align*}
 q(u,y)&=\sum_x\mu(x)
       e^{y^{\mathsf T}K^{1/2}x-u x^{\mathsf T}Kx/2},\\
 p(u,y)&=\sum_x\mu(x)f(x)
       e^{y^{\mathsf T}K^{1/2}x-u x^{\mathsf T}Kx/2}.
\end{align*}
These functions give the density of $Y_u$ as $q\phi_u$ and the
posterior expectation as $p/q=\nu_{u,y}[f]$.  Their finite exponential
sums show directly that both $p$ and $q$ solve the backward heat equation:
\[
 (\partial_u+\tfrac12\Delta_y)p=0,\qquad
 (\partial_u+\tfrac12\Delta_y)q=0.
\]
Writing $r=p/q$, the chain rule gives the exact identity
\[
 (\partial_u+\tfrac12\Delta_y)\bigl(p\log(p/q)\bigr)
     =\frac12\frac{q\|\nabla_y r\|^2}{r}.
\]
Since $\partial_u\phi_u=\tfrac12\Delta_y\phi_u$, differentiation and
integration by parts yield
\[
 \frac{d}{du}\int p\log(p/q)\,\phi_u\,dy
   =\frac12\int\frac{q\|\nabla_y r\|^2}{r}\phi_u\,dy.
\]
Differentiation and integration by parts are legitimate: the spin sums
are finite, $f$ is bounded above and below by positive constants, and
the remaining dependence on $y$ grows at most exponentially in $\|y\|$.
Gaussian tails therefore control all terms.  To identify the derivative
as the claimed covariance expression, use
\[
 H(u)=\mu[f\log f]-\int p\log(p/q)\phi_u\,dy,
 \qquad \nabla_y r=K^{1/2}\Cov_{\nu_{u,y}}(X,f),
\]
which proves Equation~\eqref{rcut:eq:posterior-entropy-derivative}.
At the initial endpoint $u\downarrow0$, the observations tend to zero
in probability and the posterior tends to $\mu$.  Since the entropy
functional is bounded on this finite space,
$H(u)\to\Ent_\mu(f)$, giving the initial value for the comparison.

\medskip
\noindent\emph{Covariance controlled by entropy.}
We next bound the derivative by the expected entropy itself.
Let $\nu$ be any posterior, take a unit vector $a$, and set
$Z=a^{\mathsf T}(X-\nu[X])$.  The second derivative of the log
moment-generating function of $Z$ is the variance of $a^{\mathsf T}X$
under a linear tilt of $\nu$.  Applying the uniform covariance estimate
to these tilts and integrating twice gives
\[
 \log\nu[e^{\lambda Z}]\le\frac C2\lambda^2,
 \qquad \lambda\in\R.
\]
The change of measure $\widehat\nu=f\nu/\nu[f]$ connects this
exponential-moment bound to entropy.  Apply Jensen's inequality to the
logarithm under $\widehat\nu$ to obtain
\[
 \lambda\widehat\nu[Z]
 \le\frac{\Ent_\nu(f)}{\nu[f]}+
       \log\nu[e^{\lambda Z}]
 \le\frac{\Ent_\nu(f)}{\nu[f]}+\frac C2\lambda^2.
\]
Optimize in $\lambda$ and take $a$ parallel to $\Cov_\nu(X,f)$ when
this vector is nonzero.  This yields
\begin{equation}
 \|\Cov_\nu(X,f)\|^2\le2C\nu[f]\Ent_\nu(f).
 \label{rcut:eq:covariance-entropy}
\end{equation}
Equations~\eqref{rcut:eq:posterior-entropy-derivative} and
\eqref{rcut:eq:covariance-entropy} imply $H'(u)\ge-C\|K\|H(u)$.
Integrating from zero to one gives
\begin{equation}
 \Ent_\mu(f)\le e^{C\|K\|}H(1)\le e^{3C}H(1).
 \label{rcut:eq:posterior-entropy-comparison}
\end{equation}
We have now compared the initial entropy with the entropy at the product
endpoint.  The comparison used covariance control at every interaction
scale: controlling linear tilts only at the original interaction $J$
would leave the quadratic change in
Equation~\eqref{rcut:eq:gaussian-posterior} untreated.  It remains to
express the endpoint entropy in terms of the original one-site energies.

\medskip
\noindent\emph{Removing the observation at the product endpoint.}
For a product probability measure $\pi=\bigotimes_i\pi_i$, entropy
tensorizes:
\[
 \Ent_\pi(f)\le\sum_i\pi[\Ent_{\pi_i}(f)],
\]
where in the $i$th summand the other coordinates are fixed inside the
one-coordinate entropy.  To prove tensorization, separate one coordinate
using the entropy chain rule and apply induction to the remaining
product.  For a fixed base measure, $f\mapsto\Ent(f)$ is convex; hence
averaging the separated coordinate can only decrease the remaining
one-coordinate entropies.  The required convexity follows from joint
convexity of $(a,b)\mapsto a\log(a/b)$ with $a=f(x)$ and $b=\pi[f]$.

Apply tensorization to the product posterior at $u=1$.  We must then
average under the joint law of $(X,Y_1)$, retaining the dependence between
the observation and the spins.  For this step put $F=f(X)$,
$\varphi(z)=z\log z$,
\[
 \mathcal G_i=\sigma(X_{-i}),\qquad
 \mathcal H_i=\sigma(X_{-i},Y_1),\qquad
 e_{\mathcal A}(F)=\E[\varphi(F)\mid\mathcal A]
                   -\varphi(\E[F\mid\mathcal A]).
\]
Conditional Jensen and the tower property give
\begin{align*}
 \E[e_{\mathcal H_i}(F)\mid\mathcal G_i]
 &=\E[\varphi(F)\mid\mathcal G_i]
   -\E[\varphi(\E[F\mid\mathcal H_i])\mid\mathcal G_i]\\
 &\le\E[\varphi(F)\mid\mathcal G_i]
       -\varphi(\E[F\mid\mathcal G_i])
 =e_{\mathcal G_i}(F).
\end{align*}
Thus the additional conditioning on the observation decreases the
averaged conditional entropy, even though the observation and the
complementary spins are dependent.  Combining this fact with the
posterior tensorization gives
\[
 H(1)\le\sum_i\E e_{\mathcal H_i}(F)
      \le\sum_i\E e_{\mathcal G_i}(F)
      =\sum_i\mu\bigl[\Ent_{\mu(\cdot\mid X_{-i})}(f)\bigr].
\]
We can now compare each original one-coordinate entropy with the
corresponding heat-bath energy.  Concavity of the logarithm gives
\[
 \Ent(f)\le\Cov(f,\log f),
\]
because the difference of the right and left sides is
$\E f\,[\log\E f-\E\log f]\ge0$.
Moreover, summing these conditional covariances gives
\[
 \sum_i\mu[\Cov(f,\log f\mid X_{-i})]
 =\sum_i\mu[f\mathsf D_i\log f]
 =\mu[f(-L)\log f].
\]
The product-endpoint estimate is therefore bounded by the entropy
dissipation of the original generator.  Substituting into
Equation~\eqref{rcut:eq:posterior-entropy-comparison} proves
Equation~\eqref{rcut:eq:entropy-inequality}, with the permissible choice
$C_1=e^{3/\kappa_\beta}$.

\medskip
\noindent\emph{Entropy decay and mixing.}
For a point start, it remains to bound the initial relative entropy
$\log(1/\mu(x))$ uniformly in the starting point.  The identity
$\|x\|^2=n$ gives
\[
 \left|\frac12x^{\mathsf T}Jx\right|\le\frac n2\|J\|,
 \qquad
 \log\frac1{\mu(x)}\le n(\log2+\|J\|)\le n(\log2+1).
\]
Reversibility lets us evolve the density $f_t=d\rho_{t,x}/d\mu$ using
the same semigroup $S_t$.  For $t>0$ the density is strictly positive,
so the entropy inequality applies to its derivative:
\[
 \frac{d}{dt}\Ent_\mu(f_t)
   =-\mu[f_t(-L)\log f_t]
   \le-C_1^{-1}\Ent_\mu(f_t).
\]
Integrating this differential inequality and using continuity at zero
establishes Equation~\eqref{rcut:eq:entropy-decay}, with
$C_2=\log2+1$.

To finish, coarsen the measure to an event and its complement; relative
entropy then dominates the resulting binary relative entropy.  For
fixed $q\in(0,1)$, this binary entropy, viewed as a function of $p$,
has value and derivative zero at $p=q$ and second derivative
$1/[p(1-p)]\ge4$.  It is consequently at least $2(p-q)^2$.
Choosing an event that attains total variation and using the entropy
decay just proved yields
\[
 d_c(t)\le\sqrt{\frac{C_2n}{2}}\,e^{-t/(2C_1)},
\]
which proves the stated mixing upper bound.
\end{proof}
